\documentclass[aspectratio=169,12pt]{beamer}
\usepackage{amsmath}
\usepackage{amssymb}
\usepackage{array}
\usepackage[most]{tcolorbox}
\usepackage{graphicx}
\usepackage[sfdefault,lf]{FiraSans}
\renewcommand{\familydefault}{\sfdefault}
\setlength{\parindent}{0pt}
\everymath{\displaystyle}

\setbeamertemplate{navigation symbols}{}
\setbeamertemplate{footline}{}
\setbeamertemplate{headline}{}
\setbeamertemplate{blocks}[rounded][shadow=false]

\definecolor{mmpageWhite}{HTML}{FCFBF7}
\definecolor{mmtanSoft}{HTML}{F7F5EF}
\definecolor{mmgreenDeep}{HTML}{244B2C}
\definecolor{mmgreenLeaf}{HTML}{5D8A56}
\definecolor{mmgreenSoft}{HTML}{E4EFD9}
\definecolor{mmtext}{HTML}{163221}
\definecolor{mmwarn}{HTML}{8F2F36}
\definecolor{mmwarnSoft}{HTML}{F8E8EA}

\setbeamercolor{background canvas}{bg=mmpageWhite}
\setbeamercolor{normal text}{fg=mmtext,bg=mmpageWhite}
\setbeamercolor{block title}{fg=mmgreenDeep,bg=mmgreenSoft}
\setbeamercolor{block body}{fg=mmtext,bg=mmtanSoft}
\setbeamercolor{block title example}{fg=mmgreenDeep,bg=mmgreenSoft}
\setbeamercolor{block body example}{fg=mmtext,bg=white}
\setbeamercolor{block title alerted}{fg=mmwarn,bg=mmwarnSoft}
\setbeamercolor{block body alerted}{fg=mmtext,bg=white}
\setbeamercolor{itemize item}{fg=mmgreenDeep}
\setbeamercolor{itemize subitem}{fg=mmgreenDeep}

\newtcolorbox{MMCard}[2][]{enhanced,breakable=false,colback=#2,colframe=black,boxrule=1.5pt,arc=8pt,left=5pt,right=5pt,top=4pt,bottom=4pt,before skip=0pt,after skip=0pt,#1}
\newcommand{\KoruIcon}{\raisebox{-0.2em}{\includegraphics[height=1.05em]{../../../manamaths/SITE/header-logo.png}}}
\newcommand{\NotesTitle}[1]{{\colorbox{mmtanSoft}{\parbox{0.975\linewidth}{\vspace{0.14em}\hspace{0.25em}{\LARGE\bfseries\textcolor{mmgreenDeep}{#1}\hfill\KoruIcon}\vspace{0.14em}}}\par\vspace{0.18em}{\color{mmgreenLeaf}\rule{\linewidth}{2.0pt}}\par\vspace{0.12em}}}
\newcommand{\SectionTitle}[1]{{\large\bfseries\textcolor{mmgreenDeep}{#1}}\par\vspace{0.04em}}
\newcommand{\GreenDivider}{\par\vspace{0.01em}{\color{mmgreenLeaf}\rule{\linewidth}{2.4pt}}\par\vspace{0.03em}}
\newcommand{\KeyIdea}[1]{\begin{MMCard}[title={\textbf{Key idea}},colbacktitle=mmgreenSoft,coltitle=mmgreenDeep]{mmtanSoft}\small #1\end{MMCard}}
\newcommand{\WorkedExample}[2]{\begin{MMCard}[title={\textbf{#1}},colbacktitle=mmgreenSoft,coltitle=mmgreenDeep]{white}\small #2\end{MMCard}}
\newcommand{\CommonMistake}[1]{\begin{MMCard}[title={\textbf{Common mistake}},colbacktitle=mmwarnSoft,coltitle=mmwarn]{white}\small #1\end{MMCard}}
\newcommand{\TryThese}[1]{\begin{MMCard}[title={\textbf{Try these}},colbacktitle=mmgreenSoft,coltitle=mmgreenDeep]{mmtanSoft}\small #1\end{MMCard}}
\newcommand{\StepLine}[2]{\textbf{#1.} #2\par\vspace{0.03em}}

\usepackage{tikz}

\setbeamertemplate{background}{
 \begin{tikzpicture}[remember picture,overlay]
 \draw[
 black,
 line width=1.5pt,
 rounded corners=10pt
 ]
 ([xshift=0.45cm,yshift=-0.45cm]current page.north west)
 rectangle
 ([xshift=-0.45cm,yshift=0.45cm]current page.south east);
 \end{tikzpicture}
}

\begin{document}

\begin{frame}[t,shrink=8]
\NotesTitle{Notes \& Steps}
\footnotesize
\begin{columns}[T,onlytextwidth]
\begin{column}{0.48\textwidth}
\KeyIdea{A power (or exponent) tells you how many times to multiply a number by itself. A root asks: what number, multiplied by itself, gives the original?}
\SectionTitle{Steps — evaluate powers}
\StepLine{1}{Identify the base (the big number) and the exponent (the small number).}
\StepLine{2}{Multiply the base by itself exponent times.}
\StepLine{3}{Example: \(2^{3} = 2 \times 2 \times 2 = 8\).}
\SectionTitle{Steps — square roots}
\StepLine{1}{Ask: what number times itself equals the given?}
\StepLine{2}{Example: \(\sqrt{25} = 5\) because \(5 \times 5 = 25\).}
\end{column}
\begin{column}{0.48\textwidth}
\SectionTitle{Key facts}
\begin{itemize}
  \setlength{\itemsep}{0.12em}
  \item \(2 \times 2 \times 2 \times 2 = 2^{4}\)
  \item \(2^{3} = 8\)
  \item \(3^{2} = 9\)
  \item \(\sqrt{25} = 5\)
  \item \(4^{1} = 4\) (any number to the power 1 is itself)
  \item \(\sqrt{4} = 2\) (rational — an integer)
\end{itemize}
\end{column}
\end{columns}
\GreenDivider
\vfill
\begin{columns}[b,onlytextwidth]
\begin{column}{0.48\textwidth}
\CommonMistake{Thinking \(2^{3} = 6\) (multiplying base by exponent). Remember: \(2^{3} = 2 \times 2 \times 2 = 8\), not \(2 \times 3 = 6\).}
\end{column}
\begin{column}{0.48\textwidth}
\TryThese{\begin{enumerate}
  \item Write \(2 \times 2 \times 2 \times 2\) using a power.
  \item What is \(2^{3}\)?
  \item What is \(\sqrt{25}\)?
\end{enumerate}}
\end{column}
\end{columns}
\end{frame}

\begin{frame}[t,shrink=12]
\NotesTitle{Notes \& Steps}
\begin{columns}[T,onlytextwidth]
\begin{column}{0.48\textwidth}
\WorkedExample{Example 1: square numbers}{Evaluate \(3^{2}\).\[3^{2} = 3 \times 3 = 9\]Answer: \(9\). \(9\) is called a square number.}
\end{column}
\begin{column}{0.48\textwidth}
\WorkedExample{Example 2: cube numbers}{Calculate \(5^{3}\).\[5^{3} = 5 \times 5 \times 5 = 125\]Answer: \(125\). \(125\) is called a cube number.}
\end{column}
\end{columns}
\GreenDivider
\TryThese{\begin{enumerate}
  \item Evaluate \(3^{2}\).
  \item Find \(\sqrt{36}\).
  \item Name one cube number.
\end{enumerate}}
\end{frame}

\begin{frame}[t,shrink=12]
\NotesTitle{Notes \& Steps}
\begin{columns}[T,onlytextwidth]
\begin{column}{0.48\textwidth}
\WorkedExample{Example 3: powers of 10}{Calculate \(10^{3}\).\[10^{3} = 10 \times 10 \times 10 = 1000\]Pattern: the exponent tells how many zeros.}
\end{column}
\begin{column}{0.48\textwidth}
\WorkedExample{Example 4: expressing as a power}{Write \(81\) as a power of \(3\).\[81 = 3 \times 3 \times 3 \times 3 = 3^{4}\]Answer: \(3^{4}\).}
\end{column}
\end{columns}
\GreenDivider
\CommonMistake{Confusing \(-3^{2}\) with \((-3)^{2}\). \(-3^{2} = -9\) (the square happens first). \((-3)^{2} = 9\) (the whole negative is squared).}
\end{frame}

\end{document}
